% ---------------------------------------------------------------- remote shell in detail
\subsection{The encrypted remote shell}
\label{sec:shellimpl}

\textbf{Server.} \texttt{server::run} loads (or creates) the host key, listens
on TCP port 2222 and starts one thread per accepted connection. That thread runs
the handshake and then the request loop of \texttt{session.rs}: \texttt{Upload}
and \texttt{Download} are answered and the loop continues, \texttt{Disconnect}
ends it, and \texttt{Shell} hands the connection to \texttt{shell::run}. The
shell starts in the folder where the server was started, with all three
standard streams connected to pipes. The connection is then split into its
\texttt{Sender} and \texttt{Receiver} halves (Section~\ref{sec:framing}), so
that reading and writing can happen in different threads, each half with its
own cipher.

\textbf{Client.} \texttt{rsh connect} loads \texttt{id\_rsh}, connects, runs the
handshake, sends \texttt{Shell} and waits for \texttt{Success}. It then also
splits the connection: one thread reads the local keyboard and sends
\texttt{Data}, while the main thread receives messages and writes \texttt{Data}
to the screen until \texttt{Exit} arrives. Figure~\ref{fig:threads} shows all
the threads of one session.

\begin{figure}[H]
\centering
\resizebox{\linewidth}{!}{%
\begin{tikzpicture}[
  font=\footnotesize, >={Stealth[length=2mm]},
  th/.style={draw, rounded corners, fill=blue!7, align=center, minimum width=2.7cm, minimum height=0.8cm},
  io/.style={draw, align=center, minimum width=1.5cm, minimum height=0.7cm, fill=gray!10},
  lock/.style={draw, dashed, align=center, minimum width=1.4cm, minimum height=0.6cm},
]
  % client
  \node[font=\bfseries\small] at (0.6,2.1) {client};
  \node[io] (kbd) at (-1.4,1.2) {keyboard};
  \node[io] (scr) at (-1.4,-1.2) {screen};
  \node[th] (fwd) at (1.4,1.2) {stdin forwarder\\(thread)};
  \node[th] (cmain) at (1.4,-1.2) {main thread\\(receiver)};
  \draw[->] (kbd) -- (fwd);
  \draw[->] (cmain) -- (scr);

  % wire
  \draw[->, very thick, blue!60!black] (2.75,1.2) -- node[above, align=center]{\texttt{Data}, \texttt{Eof}, \texttt{Disconnect}\\[-1pt] c2s keystream} (6.55,1.2);
  \draw[<-, very thick, blue!60!black] (2.75,-1.2) -- node[below, align=center]{s2c keystream\\[-1pt] \texttt{Data}, \texttt{Exit}} (6.55,-1.2);
  \node[gray, align=center] at (4.65,0) {one TCP connection\\(encrypted frames)};

  % server
  \node[font=\bfseries\small] at (8.7,2.1) {server};
  \node[th] (sess) at (7.9,1.2) {session thread\\(receiver)};
  \node[lock] (snd) at (7.9,-1.2) {\texttt{Sender}\\ (mutex)};
  \node[th] (pout) at (10.6,-0.4) {stdout pump\\+ echo filter};
  \node[th] (perr) at (10.6,-1.55) {stderr pump};
  \node[th] (watch) at (10.6,-2.7) {exit watcher};
  \node[io, minimum height=4.5cm, minimum width=1.3cm] (sh) at (13.6,-0.6) {shell\\process};
  \draw[->] (sess) -- node[above]{stdin pipe} (sh.north west |- sess);
  \draw[<-] (pout) -- node[above, pos=0.55]{stdout} (pout -| sh.west);
  \draw[<-] (perr) -- node[above, pos=0.55]{stderr} (perr -| sh.west);
  \draw[->] (pout.west) -- (snd.east);
  \draw[->] (perr.west) -- (snd.east);
  \draw[->] (watch.west) -| (snd.south);
  \draw[->, dotted] (watch.east) -| node[pos=0.25, below]{\texttt{wait()}} (sh.south);
  \draw[->] (snd.west) -- (6.55,-1.2);
  \draw[<-] (sess.west) -- (6.55,1.2);
\end{tikzpicture}}
\caption{Threads of one interactive session. The three server threads that send share one \texttt{Sender} behind a mutex, so their messages never interleave inside a frame and the keystream positions stay in order.}
\label{fig:threads}
\end{figure}

The mutex around the \texttt{Sender} matters for the encryption: the cipher
state moves forward with every byte, so two threads must never encrypt at the
same time. Each frame is encrypted and written while the lock is held, which
makes the order of frames on the wire the same as the order of keystream use.

\Needspace{18\baselineskip}
\begin{lstlisting}[caption={The server's input loop (\texttt{server/shell.rs}).}]
let result = loop {
    match receiver.recv() {                         // decrypt with the c2s keystream
        Ok(Message::Data(data)) => {
            if let Some(input) = stdin.as_mut() {
                echo.lock().unwrap().expect(&data); // Windows: note the echo
                if input.write_all(&data).and_then(|_| input.flush()).is_err() {
                    stdin = None;                   // shell closed its input
                }
            }
        }
        Ok(Message::Eof) => stdin = None,           // dropping the pipe sends EOF
        Ok(Message::Disconnect) => break Ok(()),
        Ok(_) => break Err(invalid("unexpected message during shell")),
        Err(e) => break Err(e),                     // connection lost
    }
};
drop(stdin);
let _ = child.lock().unwrap().kill();               // never leave a shell behind
\end{lstlisting}

\textbf{How a session ends.} Every way out ends with the shell gone and the
client told why:

\begin{table}[H]
\centering
\footnotesize
\begin{tabularx}{\textwidth}{@{}l X@{}}
\toprule
Event & What happens \\
\midrule
User types \texttt{exit 3} & the shell exits; its pipes close, both pumps finish, the watcher collects code 3 and sends \texttt{Exit(3)}; the client sends \texttt{Disconnect} and exits with code 3 \\
End of input on the client (Ctrl+D, end of a pipe) & the forwarder sends \texttt{Eof}; the server drops the shell's stdin pipe; \texttt{sh} sees end of input and exits as above \\
Client killed or network lost & the server's \texttt{recv} fails, the loop ends and the server kills the shell \\
Shell killed on the server & it has no exit code; the watcher sends \texttt{Exit(-1)} \\
Shell cannot be started & the server sends \texttt{Failure("cannot start shell: \ldots")} instead of \texttt{Success} \\
\bottomrule
\end{tabularx}
\caption{Ending an interactive session.}
\end{table}

\textbf{Echo filter.} With no terminal attached, \texttt{cmd.exe} echoes every
command line it reads to its output, so the user would see each command twice.
On Windows the server therefore records each complete input line, for example
\texttt{dir\textbackslash r\textbackslash n}, and the stdout pump removes
exactly that byte sequence when it arrives. As soon as the output differs from
what was expected, the held bytes are released, so real output is never lost.
On other systems the filter does nothing, because \texttt{sh} does not echo.

\textbf{File copy.} \texttt{rcp} uses the same connection and handshake but
sends one \texttt{Upload} (path, file name, contents) or \texttt{Download}
(path) and waits for \texttt{Success}, \texttt{File} or \texttt{Failure},
then sends \texttt{Disconnect}. The whole file travels in one encrypted
message, so its contents are protected exactly like the shell traffic.
